% © 2026 Ing. David Jaroš — CC BY-NC-ND 4.0
%
% This work is licensed under a Creative Commons Attribution-NonCommercial-NoDerivatives
% 4.0 International License (CC BY-NC-ND 4.0).
%
% See LICENSE.md for full license text.
%
% Classification: sections 1-2 [STD+MC], sections 3-4 [MC/SPECULATIVE].

\documentclass[12pt,a4paper]{article}
\usepackage[a4paper, margin=2.5cm]{geometry}
\usepackage{amsmath,amssymb}
\usepackage{hyperref}
\title{\textbf{NCG spektrální trojice pro UBT}}
\author{Ing.\ David Jaro\v{s}}
\date{2026-05-11}

\begin{document}
\maketitle

\section{Definice spektrální trojice z UBT \,[STD+MC]}

Navrhujeme Connesovu spektrální trojici $(A,H,D)$:
\begin{align}
A &= \mathbb C\otimes\mathbb H \cong \mathrm{Mat}(2,\mathbb C),\\
H &= L^2(M,S\otimes(\mathbb C\otimes\mathbb H)),\\
D &= \text{Diracův operátor na }M\text{ s hodnotami v }\mathbb C\otimes\mathbb H.
\end{align}
Tento zápis je přímou NCG formulací algebraického základu UBT.

\section{Porovnání s Connes--Chamseddine \,[STD+MC]}

Connes--Chamseddine používá produktovou algebru
$C^\infty(M)\otimes(\mathbb C\oplus\mathbb H\oplus M_3(\mathbb C))$.
UBT naopak používá jednoduchou algebru
$\mathbb C\otimes\mathbb H\cong \mathrm{Mat}(2,\mathbb C)$.

Pracovní strukturální teze:
\begin{itemize}
  \item CC model: Standard Model je zaveden přes produktovou interní algebru.
  \item UBT: jednotný algebraický objekt dává gravitační i gauge sektor v jedné
  kostře.
\end{itemize}
Zda je tento popis fakticky jednodušší i predikčně silnější, zůstává test
pro další práci.

\section{Spektrální akce \,[MC/SPECULATIVE]}

Použijeme standardní tvar
\[
S_{\mathrm{spec}}[D]=\mathrm{Tr}\!\left[f\!\left(\frac{D}{\Lambda}\right)\right]
\sim a_0\Lambda^4+a_2\Lambda^2+a_4\ln\Lambda+\cdots.
\]
Klíčový otevřený úkol: porovnat Seeley--DeWittovy koeficienty $a_k$ s
koeficienty v UBT akci $S[\Theta]$ (gravitační a gauge část).

\section{Winding spektrum a link na $B_0$ \,[MC/SPECULATIVE]}

Pro imaginární časový kruh $S^1_\psi$ s poloměrem $R_\psi$:
\[
\lambda_n=\frac{n}{R_\psi},\qquad n\in\mathbb Z,
\]
a
\[
\mathrm{Tr}\!\left[f\!\left(\frac{D}{\Lambda}\right)\right]
\supset \sum_{n\in\mathbb Z} f\!\left(\frac{n}{R_\psi\Lambda}\right).
\]
Poissonovo přesčítání může generovat logaritmický člen, kandidáta na vazbu
ke koeficientu $B_0$ v alpha tracku. Otevřený test:
\[
\sum_n f\!\left(\frac{n}{R_\psi\Lambda}\right)\stackrel{?}{\longrightarrow}B_0\ln\Lambda.
\]
Pokud se potvrdí, vznikne přímý most mezi NCG spektrální akcí a
vacuum-polarizačním koeficientem v T3\_ALPHA.

\end{document}

